% =============================================================
%  附录 D  记录表
% =============================================================
\Chapter{附录 D\quad 记录表和速查卡}

\subsection*{D1\quad 走完全册打钩表}

\noindent 上完一次就打个钩，再画颗星。\textbf{41 次全打满，这一册就过完了。}

\begin{center}
\small
\renewcommand{\arraystretch}{1.45}
\begin{tabular}{@{}p{4.4cm}|*{7}{>{\centering\arraybackslash}p{1.15cm}|}@{}}
\hline
\textbf{哪一章} & \textbf{1} & \textbf{2} & \textbf{3} & \textbf{4} & \textbf{5} &
\textbf{6} & \textbf{7} \\
\hline
第一单元\ 加与减（二） & & & & & & & \\ \hline
第二单元\ 测量（一） & & & & & & & \\ \hline
第三单元\ 数一数与乘法 & & & & & & & \\ \hline
第四单元\ 乘法口诀（一） & & & & & & & \\ \hline
第五单元\ 分一分与除法 & & & & & & & \\ \hline
第六单元\ 图形的运动（一） & & & & & & & \\ \hline
第七单元\ 乘法口诀（二） & & & & & & & \\ \hline
第八单元\ 乘除法的应用（一） & & & & & & & \\ \hline
两次综合实践 & & & & & & & \\ \hline
期末复习 & & & & & & & \\ \hline
\end{tabular}
\end{center}

\vspace{4pt}
\noindent{\small 第六单元只有 3 课，第八单元 4 课，综合实践 2 课，
其余 5--6 课。空格多的地方不用硬填满。}

\subsection*{D2\quad 错题本：一次只写一句话}

\noindent\textbf{规矩就一条：只写错在哪，不抄整道题。}
同一个错因出现三次，就停下来回去补那一课。

\begin{center}
\small
\renewcommand{\arraystretch}{1.9}
\begin{tabular}{@{}p{1.7cm}|p{6.6cm}|p{1.9cm}|p{1.4cm}@{}}
\hline
\textbf{日期} & \textbf{错在哪一步（一句话）} & \textbf{重做日期} & \textbf{会了吗} \\
\hline
 & & & \\ \hline
 & & & \\ \hline
 & & & \\ \hline
 & & & \\ \hline
 & & & \\ \hline
 & & & \\ \hline
 & & & \\ \hline
 & & & \\ \hline
 & & & \\ \hline
 & & & \\ \hline
 & & & \\ \hline
 & & & \\ \hline
\end{tabular}
\end{center}

\vspace{2pt}
\noindent\textbf{好的写成长这样}：
\begin{itemize}
  \item “48+25，个位 8+5=13，把进的 1 忘了加到十位”
  \item “52-37，十位 5 借走了 1，还是按 5 减的”
  \item “把‘是 3 倍’当成了‘多 3 个’”
\end{itemize}
\noindent\textbf{没用的写成长这样}：“粗心”“算错了”“不认真”——一条都别写，
写了等于没写。

\clearpage
\subsection*{D3\quad 星星表（攒够 20 颗换奖励）}

\begin{center}
\renewcommand{\arraystretch}{2.5}
\begin{tabular}{|*{5}{>{\centering\arraybackslash}p{2.7cm}|}}
\hline
 & & & & \\ \hline
 & & & & \\ \hline
 & & & & \\ \hline
 & & & & \\ \hline
\end{tabular}
\end{center}

\vspace{6pt}
\begin{center}
\begin{tabular}{@{}p{9cm}p{6cm}@{}}
\toprule
\textbf{怎么得} & \textbf{几颗} \\
\midrule
一次课顺顺当当上完（不哭、不跑、不超时） & \starx \\
一次“5 分钟快检”全对 & \starx \\
\textbf{小豆角自己发现并改了一个错} & \starx\starx \\
一周每天都做了 3 分钟口算卡 & \starx \\
\bottomrule
\end{tabular}
\end{center}

\clearpage
\subsection*{D4\quad 剪下来贴冰箱的那张}

\begin{center}
\begin{tcolorbox}[enhanced,colback=white,colframe=mainc,boxrule=1.2pt,
                  arc=2mm,width=16.6cm,left=10pt,right=10pt,top=8pt,bottom=8pt]
\begin{center}
{\heiti\large 每次二十分钟，就这么切}\\[4pt]
\renewcommand{\arraystretch}{1.35}
\begin{tabular}{@{}p{2.2cm}p{2.6cm}p{9.6cm}@{}}
\toprule
3 分钟 & 抽卡口算 & 抽 10 张，只记对错，当场不纠 \\
2 分钟 & 小豆角自己读题 & 只说一句“把数字圈出来”，豆角爸爸不念题 \\
8 分钟 & \textbf{上手摆} & \textbf{全场最重要，豆角爸爸不可以说话} \\
5 分钟 & 落到纸上 & 写算式／竖式／口诀，豆角爸爸不代笔 \\
2 分钟 & 收尾 & 三问：学了啥／哪步易错／还来一次吗 \\
\bottomrule
\end{tabular}
\end{center}
\end{tcolorbox}
\end{center}

\begin{center}
\begin{tcolorbox}[enhanced,colback=white,colframe=pitc,boxrule=1.2pt,
                  arc=2mm,width=16.6cm,left=10pt,right=10pt,top=8pt,bottom=8pt]
\begin{center}
{\heiti\large 六件豆角爸爸最不该做的事}\\[4pt]
\begin{tabular}{@{}p{0.6cm}p{14.2cm}@{}}
1 & 把规则直接说出口（让小豆角自己摆出来） \\
2 & 拿口诀顶替理解（每句新口诀先圈图） \\
3 & 管小豆角掰手指（二年级上学期正常，别管） \\
4 & 用题量顶思考（一种类型做透 5 道） \\
5 & 边做边纠错（整道题做完再回头看） \\
6 & 盯着分数（能说清“为什么”才是目标） \\
\end{tabular}
\end{center}
\end{tcolorbox}
\end{center}

\begin{center}
\begin{tcolorbox}[enhanced,colback=white,colframe=passc,boxrule=1.2pt,
                  arc=2mm,width=16.6cm,left=10pt,right=10pt,top=8pt,bottom=8pt]
\begin{center}
{\heiti\large 小豆角卡住了，就问这三句}\\[4pt]
{\kaishu ① “小豆角是怎么想的？”\quad ② “这一捆／这一份是从哪来的？”\quad
③ “小豆角能画出来给爸爸看吗？”}\\[3pt]
{\small 三句都不管用 $\to$ 换材料（小棒换豆子、纸换白板）$\to$ 还不行就今天先停，
明天再来。\textbf{停下来永远比硬推有用。}}
\end{center}
\end{tcolorbox}
\end{center}
